% Coverage: Chapter 8, Section 8.5, physical PDF pp. 376--378
% (printed pp. 353--355; boundary pages contain adjacent-section material).
% Semantic range: Section 8.5 heading through "and the Coulomb interaction is
% dropped." Inventory: equations (8.5.1)--(8.5.9), one ordinary numbered
% footnote, no figures, tables, problems, references, or unresolved uncertainties.

\subsection{The Photon Propagator}

The general Feynman rules described in Chapter 6 dictate that an internal
photon line in a Feynman diagram contributes a factor to the corresponding
term in the $S$-matrix, given by the propagator:
\begin{equation}
-\ii\Delta_{\mu\nu}(x-y)
\equiv\bra{\Omega}T\{a_\mu(x),a_\nu(y)\}\ket{\Omega},
\tag{8.5.1}\label{eq:8.5.1}
\end{equation}
where $T$ as usual denotes a time-ordered product. Inserting our formula
\eqref{eq:8.4.15} for the electromagnetic potential then yields
\begin{align}
-\ii\Delta_{\mu\nu}(x-y)
=\int\frac{\dd^3p}{(2\pi)^3 2\lvert\mathbf p\rvert}
P_{\mu\nu}(\mathbf p)\Bigl[
e^{\ii p\cdot(x-y)}\theta(x^0-y^0)
+e^{\ii p\cdot(y-x)}\theta(y^0-x^0)\Bigr],
\tag{8.5.2}\label{eq:8.5.2}
\end{align}
where
\begin{equation}
P_{\mu\nu}(\mathbf p)
\equiv\sum_{\sigma=\pm1}
e_\mu(\mathbf p,\sigma)e_\nu(\mathbf p,\sigma)^*,
\tag{8.5.3}\label{eq:8.5.3}
\end{equation}
and $p^\mu$ in the exponentials is taken with
$p^0=\lvert\mathbf p\rvert$. We recall from Eqs.~\eqref{eq:8.4.18} and \eqref{eq:8.4.17} that
\begin{align}
P_{ij}(\mathbf p)&=\delta_{ij}-\frac{p^i p^j}{\lvert\mathbf p\rvert^2},
\notag\\
P_{0\ii}(\mathbf p)&=P_{i0}(\mathbf p)=P_{00}(\mathbf p)=0.
\tag{8.5.4}\label{eq:8.5.4}
\end{align}
As we saw in Chapter 6, the theta functions in Eq.~\eqref{eq:8.5.2} may be expressed as
integrals over an independent time-component $q^0$ of an off-shell
four-momentum $q^\mu$, so that Eq.~\eqref{eq:8.5.2} may be rewritten
\begin{equation}
\Delta_{\mu\nu}(x-y)=(2\pi)^{-4}\int \dd^4q\,
\frac{P_{\mu\nu}(\mathbf q)}{q^2-\ii\epsilon}
e^{\ii q\cdot(x-y)}.
\tag{8.5.5}\label{eq:8.5.5}
\end{equation}
Thus in using the Feynman rules in momentum space, the contribution of an
internal photon line carrying four-momentum $q$ that runs between vertices
where the photon is created and destroyed by fields $a^\mu$ and $a^\nu$ is
\begin{equation}
\frac{-\ii}{(2\pi)^4}\frac{P_{\mu\nu}(\mathbf q)}{q^2-\ii\epsilon}.
\tag{8.5.6}\label{eq:8.5.6}
\end{equation}

It will be very useful (though apparently perverse) to rewrite Eq.~\eqref{eq:8.5.4} as
\begin{equation}
P_{\mu\nu}(\mathbf q)=\eta_{\mu\nu}
+\frac{q^0q_\mu n_\nu+q^0q_\nu n_\mu-q_\mu q_\nu+q^2n_\mu n_\nu}
       {\lvert\mathbf q\rvert^2},
\tag{8.5.7}\label{eq:8.5.7}
\end{equation}
where $n^\mu=(1,0,0,0)$ is a fixed time-like vector, $q^2$ as usual is
$\lvert\mathbf q\rvert^2-(q^0)^2$, but $q^0$ is here entirely
arbitrary. We shall choose $q^0$ in Eq.~\eqref{eq:8.5.7} to be given by
four-momentum conservation: it is the difference of the matter $p^0$s flowing
in and out of the vertex where the photon line is created. The terms
proportional to $q_\mu$ and/or $q_\nu$ then do not contribute to the $S$-matrix,
because the factors $q_\mu$ or $q_\nu$ act like derivatives $\partial_\mu$ and
$\partial_\nu$, and the photon fields $a_\mu$ and $a_\nu$ are coupled to
currents $j^\mu$ and $j^\nu$ that satisfy the conservation condition
$\partial_\mu j^\mu=0$.
\footnote{This argument as given here is little better than hand-waving. The
result has been justified by a detailed analysis of Feynman diagrams~\hyperref[ref:8.3]{[3]}, but
the easiest way to treat this problem is by path-integral methods, as discussed
in Section 9.6.}
The term proportional to $n_\mu n_\nu$ contains a factor $q^2$ that cancels the
$q^2$ in the denominator of the propagator, yielding a term that is the same
as would be produced by a term in the action:
\[
-\ii\frac12\int \dd^4x\int \dd^4y\,
[-ij^0(x)][-ij^0(y)]\frac{-\ii}{(2\pi)^4}
\int\frac{\dd^4q}{\lvert\mathbf q\rvert^2}e^{\ii q\cdot(x-y)}.
\]
The integral over $q^0$ here yields a delta function in time, so this is
equivalent to a correction to the interaction Hamiltonian $H_I(t)$, of the
form
\[
-\frac12\int \dd^3x\int \dd^3y\,
\frac{j^0(\mathbf x,t)j^0(\mathbf y,t)}
     {4\pi\lvert\mathbf x-\mathbf y\rvert}.
\]
This is just right to cancel the Coulomb interaction \eqref{eq:8.4.25}. Our result is
that the photon propagator can be taken effectively as the covariant quantity
\begin{equation}
\Delta^{\mathrm{eff}}_{\mu\nu}(x-y)
=(2\pi)^{-4}\int \dd^4q\,
\frac{\eta_{\mu\nu}}{q^2-\ii\epsilon}e^{\ii q\cdot(x-y)}
\tag{8.5.8}\label{eq:8.5.8}
\end{equation}
with the Coulomb interaction dropped from now on. We see that the apparent
violation of Lorentz invariance in the instantaneous Coulomb interaction is
cancelled by another apparent violation of Lorentz invariance, that as noted
in Section 5.9 the fields $a^\mu(x)$ are not four-vectors, and therefore have
a non-covariant propagator. From a practical point of view, the important
point is that in the momentum space Feynman rules, the contribution of an
internal photon line is simply given by
\begin{equation}
\frac{-\ii}{(2\pi)^4}\frac{\eta_{\mu\nu}}{q^2-\ii\epsilon}
\tag{8.5.9}\label{eq:8.5.9}
\end{equation}
and the Coulomb interaction is dropped.
